\section{Explicit polylogarithmic and Gamma correlation identities}
\label{tc:section}

The preceding transport theorem is an all-index result. This section
makes several consequences explicit, including an ordinary convergent
Gamma integral and a root-of-unity example whose lower terms reduce to
classical Stieltjes constants and logarithms.

\subsection{Finite arithmetic reduction of weighted Tornheim kernels}
The weighted double kernel that appears in a direct Fourier treatment
is itself a finite transform of \eqref{tr:Tdef}. Let $p,q\geq1$ be
integers and define
\begin{equation}\label{tc:weighted}
 \mathcal T_{p,q}(A,B,C;X,Y)
   =\sum_{m,n\geq1}\frac{X^mY^n}{m^An^B(pm+qn)^C}
\end{equation}
in an absolutely convergent region. Fix any roots $\xi^p=X$ and
$\eta^q=Y$.

\begin{proposition}[A finite root filter]\label{tc:root-filter}
If $|X|=|Y|=1$ and $X,Y\ne1$, then
\begin{align}\label{tc:weighted-reduction}
 \mathcal T_{p,q}(A,B,C;X,Y)
 =p^{A-1}q^{B-1}
   \sum_{h=0}^{p-1}\sum_{k=0}^{q-1}
     \mathcal T(A,B,C;\xi e(h/p),\eta e(k/q)).
\end{align}
The expression is independent of the chosen roots and is entire in
$(A,B,C)$. In particular, for $N\geq0$,
\begin{align}\label{tc:weighted-negative}
 \mathcal T_{p,q}(A,B,-N;X,Y)
 =\sum_{j=0}^N\binom Nj p^jq^{N-j}
             \Li_{A-j}(X)\Li_{B-N+j}(Y).
\end{align}
All $A$- and $B$-derivatives of this last identity are valid.
\end{proposition}
\begin{proof}
Substitute $M=pm$ and $N'=qn$ in \eqref{tc:weighted}. The indicator
of $p\mid M$ is $p^{-1}\sum_he(hM/p)$, and similarly for $q\mid N'$.
The factors $m^{-A}=p^AM^{-A}$ and $n^{-B}=q^B(N')^{-B}$ give
\eqref{tc:weighted-reduction}. Changing a chosen root permutes the
finite sum. No new color equals one because $X,Y\ne1$; hence
Lemma~\ref{tr:T-entire} proves entire continuation. Expanding
$(pm+qn)^N$ in the original convergent region proves
\eqref{tc:weighted-negative}, and continuation gives the full claim.
\end{proof}

Finite congruence conditions on the two indices can be imposed by a
further finite root filter. This explains why integral weights and
congruence restrictions arising from the unequal-frequency lattice can
be handled without introducing a new family of transcendental kernels.
It does not assert that the finite list produced in this way has no
further relations.

The derivative in the \emph{third} order requires separate data.
Differentiating an identity known only at $C=-N$ is not legitimate.
An exact formula specifying that additional data follows from the
same Mellin continuation. Put
$f(t)=\Li_A(Xe^{-t})\Li_B(Ye^{-t})$ and $f_j=[t^j]f(t)$. Then
\begin{align}\label{tc:normal-derivative}
 \left.\partial_C\mathcal T(A,B,C;X,Y)\right|_{C=-N}
 =(-1)^NN!\Bigg\{&
  \int_0^1t^{-N-1}\left(f(t)-\sum_{j=0}^Nf_jt^j\right)\dd t
  +\int_1^\infty t^{-N-1}f(t)\dd t\notag\\
 &+\sum_{j=0}^{N-1}\frac{f_j}{j-N}
        -\psi(N+1)f_N\Bigg\}.
\end{align}
Both integrals converge ordinarily. To prove this, use the Taylor
subtraction in Lemma~\ref{tr:T-entire} and
\[
 \frac1{\Gamma(-N+v)}
   =(-1)^NN!v\{1-\psi(N+1)v+O(v^2)\}.
\]
The pole $f_N/v$ supplies the last term in
\eqref{tc:normal-derivative}. The formula displays exactly what remains
after the negative-order value reduces to products of polylogarithms.
It is an integral representation for the transverse derivative, not a
claim that this derivative reduces to the value at $C=-N$.

\subsection{An ordinary three-factor log-Gamma integral}
Define the centered log-Gamma function and a constant by
\begin{equation}\label{tc:centered-Gamma}
 \Lambda(x)=\log\Gamma(x)-\tfrac12\log(2\pi)=\zeta'(0,x),
 \qquad \kappa=\gamma+\log(2\pi).
\end{equation}
Its periodic version is integrable to every finite power near an
endpoint. The identity below therefore requires no finite part.
For a variable $A$ and an integer $p\geq1$, use the operators
\begin{equation}\label{tc:Doperators}
 D_{+,A}^{(p)}=\partial_A+\log p-\kappa-\frac{i\pi}{2},
 \qquad
 D_{-,A}^{(p)}=\partial_A+\log p-\kappa+\frac{i\pi}{2}.
\end{equation}

\begin{theorem}[Unequal-frequency centered Gamma correlation]
\label{tc:Gamma-triple}
For the disjoint grids \eqref{tr:grids},
\begin{align}
 &\int_0^1\prod_{j=1}^3\Lambda(\{p_jx+a_j\})\dd x
  \notag\\
 &\quad=-\frac1{4\pi^3}\Im
   \sum_{h_1=0}^{p_1-1}\sum_{h_2=0}^{p_2-1}\sum_{h_3=0}^{p_3-1}
   \sum_{k=1}^3
   \left.
       D_{+,A}^{(p_i)}D_{+,B}^{(p_j)}D_{-,C}^{(p_k)}
       \mathcal T(A,B,C;X_k,Y_k)
   \right|_{A=B=C=1}.
   \label{tc:Gamma-formula}
\end{align}
Here $i<j$ are the other indices, and $X_k,Y_k$ are formed from the
lifted triple $(b_{1,h_1},b_{2,h_2},b_{3,h_3})$ as in
\eqref{tr:Sdef}. The imaginary part is taken only after specialization
to these real spectral orders.
\end{theorem}
\begin{proof}
Lerch's Gamma identity in \eqref{tc:centered-Gamma} gives
\[
 \int_0^1\prod_j\Lambda(\{p_jx+a_j\})\dd x
  =-\left.\partial_{\lambda_1}\partial_{\lambda_2}
               \partial_{\lambda_3}
          K_{\boldsymbol p}(\boldsymbol\lambda;\boldsymbol a)
           \right|_{\boldsymbol\lambda=(1,1,1)}.
\]
Differentiation under the integral is justified, near each separated
endpoint, by a majorant of the form
$t^{-\rho}(1+|\log t|^3)$ with a fixed $0<\rho<1$ in a sufficiently
small neighborhood of that spectral point. Insert
\eqref{tr:Kp} and \eqref{tr:unit-K}. At $\lambda=1$,
\[
 p^{\lambda-1}\Gamma(\lambda)(2\pi)^{-\lambda}=(2\pi)^{-1},
 \qquad
 \partial_\lambda\log\{p^{\lambda-1}\Gamma(\lambda)
                  (2\pi)^{-\lambda}\}=\log p-\kappa.
\]
The phase in a region with two positive frequencies is $-i$ at this
point. Its derivatives contribute $-i\pi/2$ in the two positive
positions and $+i\pi/2$ in the negative one, giving the operators
\eqref{tc:Doperators}. The opposite region is the complex conjugate
after evaluating at real orders. Thus its pair contributes
$2\Re(-iW)=2\Im W$. The common factor $(2\pi)^{-3}$ and the initial
minus sign give \eqref{tc:Gamma-formula}.
\end{proof}

The formula is a finite Gamma--polylogarithm identity with ordinary
integrals and entire spectral derivatives. The correlation of three
uncentered $\log\Gamma$ factors follows by expanding
$\log\Gamma=\Lambda+\tfrac12\log(2\pi)$. Since
$\int_0^1\Lambda(\{px+a\})\dd x=0$, the only additional terms are
the three ordinary centered bilinear correlations and
$\tfrac18\log^3(2\pi)$. Those bilinear correlations already have
single-polylogarithm order-derivative evaluations in the incoming
dilation report \cite{incoming-sh}.

\subsection{A twelfth-root digamma example with explicit lower terms}
For pairwise distinct $b_1,b_2,b_3$ modulo one, let
\[
 \mathcal P(b_1,b_2,b_3)
  =\FP_x\int_0^1\prod_{j=1}^3\psi(\{x+b_j\})\dd x.
\]
The unit-frequency formula \eqref{tr:triple-coeff} gives
\begin{equation}\label{tc:unit-digamma}
 \mathcal P(\boldsymbol b)
  =-2\Re\sum_{k=1}^3
    \left.
      D_{+,A}^{(1)}D_{+,B}^{(1)}D_{-,C}^{(1)}
        \mathcal T(A,B,C;X_k,Y_k)
    \right|_{A=B=C=0}.
\end{equation}
This is the already established translated three-digamma coordinate
formula; our example uses it after unequal-frequency transport.

Choose
\begin{equation}\label{tc:example-data}
 (p_1,p_2,p_3)=(1,2,3),\qquad
 (a_1,a_2,a_3)=(0,\tfrac13,\tfrac14).
\end{equation}
The lifted sets are
\[
 \mathcal B_1=\{0\},\quad
 \mathcal B_2=\{\tfrac16,\tfrac23\},\quad
 \mathcal B_3=\{\tfrac1{12},\tfrac5{12},\tfrac34\}.
\]
They are disjoint, and every color in \eqref{tc:unit-digamma} is a
twelfth root of unity distinct from one. Theorem~\ref{tr:finite-transport},
with $m_1=m_2=m_3=0$ and $\gamma_0=-\psi$, gives the exact identity
\begin{align}
 &\FP_x\int_0^1\psi(x)\psi(\{2x+\tfrac13\})
                              \psi(\{3x+\tfrac14\})\dd x\notag\\
 &\quad=\frac16\sum_{b\in\mathcal B_2}\sum_{c\in\mathcal B_3}
                     \mathcal P(0,b,c)+\mathcal L,\label{tc:example}\\
 \mathcal L
 &:=\frac{\log2}{3}\sum_{c\in\mathcal B_3}J_{00}(c)
       +\frac{\log3}{2}\sum_{b\in\mathcal B_2}J_{00}(b).
       \label{tc:example-lower}
\end{align}
The last constant reduces completely. With $u=\log2$, $v=\log3$,
and $\gamma_1=\gamma_1(1)$,
\begin{align}\label{tc:lower-polynomial}
 \mathcal L={}&2(\gamma_1-\zeta(2))(u+v)
   -\gamma(6u^2+4uv+3v^2)\notag\\
 &-7u^3-7u^2v-4uv^2-\tfrac32v^3.
\end{align}
To verify the reduction, Stieltjes multiplication gives
\begin{align*}
 \gamma_0(\tfrac14)+\gamma_0(\tfrac34)&=2\gamma+6\log2,\\
 \gamma_1(\tfrac14)+\gamma_1(\tfrac34)
      &=2\gamma_1-6\gamma\log2-7\log^22,\\
 \gamma_0(\tfrac13)+\gamma_0(\tfrac23)&=2\gamma+3\log3,\\
 \gamma_1(\tfrac13)+\gamma_1(\tfrac23)
      &=2\gamma_1-3\gamma\log3-\tfrac32\log^23.
\end{align*}
Apply the same multiplication law to the two lifted grids in
\eqref{tc:example-lower} and insert \eqref{tr:J00}. Collecting terms
gives \eqref{tc:lower-polynomial}.

Independent subtracted quadrature gives the following diagnostic values:
\[
 \FP_x\int_0^1\psi(x)\psi(\{2x+\tfrac13\})
                \psi(\{3x+\tfrac14\})\dd x
       \approx77.3418540125851504884052200711,
\]
\[
 \mathcal L\approx-23.0290680355012727975296691002.
\]
These are numerical illustrations of the proved identities. Their
normalization is the stated ambient finite part. In particular, removing
the lower terms in \eqref{tc:example-lower} would change the value by
the displayed nonzero constant.
